% Einheit 1 · Kreisumfang
\setcounter{aufgabe}{0}
\einheitenkopf*{\hypertarget{e1}{}Einheit 1 von 3 · Kreisumfang}
\zweigzeile{Hier lernst du, den Umfang eines Kreises aus Radius oder Durchmesser zu berechnen und umgekehrt · neu oder schon bekannt – je nach Buch · P10 · baut auf: Rechnen mit Kommazahlen und Runden, Umfang und Fläche unterscheiden, Formeln umstellen}

\begin{aufgabe}{Ich erkenne, ob Radius oder Durchmesser gegeben ist. Kreuze an, welche Strecke eingezeichnet oder beschrieben ist. Rechne nicht.}
\begin{minipage}[t]{4.2cm}
a)\\[1mm]
\begin{kreis}[1.1]\mittelpunkt{M}\radius{40}{3\,\text{cm}}\end{kreis}\\[1mm]
\kreuz{Radius}\\
\kreuz{Durchmesser}
\end{minipage}\hspace{0.8cm}%
\begin{minipage}[t]{4.2cm}
b)\\[1mm]
\begin{kreis}[1.1]\mittelpunkt{}\durchmesser{20}{7\,\text{cm}}\end{kreis}\\[1mm]
\kreuz{Radius}\\
\kreuz{Durchmesser}
\end{minipage}\hspace{0.8cm}%
\begin{minipage}[t]{4.2cm}
c)\\[1mm]
\begin{kreis}[1.1]\mittelpunkt{M}\radius{300}{2{,}5\,\text{cm}}\end{kreis}\\[1mm]
\kreuz{Radius}\\
\kreuz{Durchmesser}
\end{minipage}

\medskip
\setcounter{teil}{3}
\begin{teile}
\teil Ein Teller misst quer durch die Mitte von Rand zu Rand 26\,cm. \quad \kreuz{Radius}\quad\kreuz{Durchmesser}
\teil Der Zirkel ist 6\,cm weit geöffnet. \quad \kreuz{Radius}\quad\kreuz{Durchmesser}
\end{teile}
\end{aufgabe}

\begin{aufgabe}{Ich erkenne, ob nach dem Rand oder nach der Fläche gefragt ist. Kreuze an, ob man den Umfang oder den Flächeninhalt des Kreises braucht. Rechne nicht.}
\begin{teile}
\teil Wie lang wird der Zaun um ein rundes Gehege? \\ \kreuz{Umfang}\quad\kreuz{Flächeninhalt}
\teil Wie viel Rasen muss man für ein rundes Beet säen? \\ \kreuz{Umfang}\quad\kreuz{Flächeninhalt}
\teil Wie viel Spitzenborte braucht man für den Rand einer runden Tischdecke? \\ \kreuz{Umfang}\quad\kreuz{Flächeninhalt}
\teil Wie viel Tomatensoße braucht man, um eine Pizza zu bestreichen? \\ \kreuz{Umfang}\quad\kreuz{Flächeninhalt}
\teil Wie weit rollt ein Reifen bei einer Umdrehung? \\ \kreuz{Umfang}\quad\kreuz{Flächeninhalt}
\end{teile}
\end{aufgabe}

\begin{aufgabe}{Ich kann einen Kreis zeichnen und Mittelpunkt, Radius und Durchmesser einzeichnen (kennst du seit Klasse 5). Zeichne mit dem Zirkel um den Punkt M.}
\begin{teile}
\teil einen Kreis mit dem Radius $r = 2{,}5$\,cm (linke Fläche)
\teil einen Kreis mit dem Durchmesser $d = 7$\,cm (rechte Fläche)
\teil Zeichne in den Kreis aus b) einen Radius und einen Durchmesser ein und beschrifte sie mit $r$ und $d$.
\end{teile}

\zeichenflaeche{6}{6}\hspace{1cm}\zeichenflaeche{8}{8}
\end{aufgabe}

\begin{aufgabe}{Ich kann aus dem Radius den Durchmesser berechnen und umgekehrt. Berechne die fehlende Größe.}
\begin{teilezwei}
\tz $r = 6$\,cm \quad \feld[cm]{d} & \tz $d = 18$\,cm \quad \feld[cm]{r} \\
\tz $d = 11$\,cm \quad \feld[cm]{r} & \tz $r = 3{,}5$\,m \quad \feld[m]{d} \\
\tz $d = 1{,}3$\,m \quad \feld[m]{r} & \tz $d = 1$\,m \quad \feld[cm]{r} \\
\end{teilezwei}
\end{aufgabe}

\begin{aufgabe}{Ich kann herausfinden, wie oft der Durchmesser in den Umfang passt. Lena hat an runden Gegenständen den Durchmesser $d$ und mit einem Faden den Umfang $u$ gemessen. Berechne jeweils $u : d$ und runde auf zwei Stellen nach dem Komma.}
\sachtabelle{lccc}{Gegenstand & $d$ in cm & $u$ in cm & $u : d$}{a) Teelicht & 3,8 & 12,0 & \leerzelle \\ b) Dose & 7,5 & 23,5 & \leerzelle \\ c) Teller & 24,0 & 75,5 & \leerzelle \\ d) Fahrradreifen & 66 & 207 & \leerzelle}
\setcounter{teil}{4}
\begin{teile}
\teil Ergänze den Satz: Der Umfang ist ungefähr \leerfeld-mal so lang wie der Durchmesser. \\ Diese Zahl heißt $\pi$ (sprich: pi). Der Taschenrechner hat eine Taste dafür.
\end{teile}
\end{aufgabe}

\begin{aufgabe}{Ich kann den Umfang eines Kreises berechnen. Rechne mit der $\pi$-Taste und runde auf eine Stelle nach dem Komma.}
\beispiel{d = 6\,\text{cm} \quad\rightarrow\quad u &= \pi \cdot 6\,\text{cm} \approx 18{,}8\,\text{cm}}
Berechne den Umfang aus dem Durchmesser: $u = \pi \cdot d$.
\begin{teilezwei}
\tz $d = 3$\,cm \quad \feld[cm]{u} & \tz $d = 7$\,cm \quad \feld[cm]{u} \\
\tz $d = 8$\,m \quad \feld[m]{u} & \tz $d = 12$\,cm \quad \feld[cm]{u} \\
\end{teilezwei}
Ist der Radius gegeben, rechnest du $u = 2 \cdot \pi \cdot r$. Berechne den Umfang.
\begin{teilezwei}
\tz $r = 9$\,cm \quad \feld[cm]{u} & \tz $d = 2{,}6$\,cm \quad \feld[cm]{u} \\
\tz $r = 0{,}75$\,m \quad \feld[m]{u} & \\
\end{teilezwei}
\end{aufgabe}

\begin{aufgabe}{Ich kann aus dem Umfang den Durchmesser und den Radius berechnen. Runde auf eine Stelle nach dem Komma.}
\beispiel{u = 22\,\text{cm} \quad\rightarrow\quad d &= 22\,\text{cm} : \pi \approx 7{,}0\,\text{cm}}
Berechne den Durchmesser.
\begin{teilezwei}
\tz $u = 44$\,cm \quad \feld[cm]{d} & \tz $u = 100$\,m \quad \feld[m]{d} \\
\end{teilezwei}
Berechne den Radius. Runde erst am Ende.
\begin{teilezwei}
\tz $u = 50$\,cm \quad \feld[cm]{r} & \\
\end{teilezwei}
Die Figur ist ein Halbkreis mit dem Durchmesser $d = 9$\,cm.

\begin{kreis}[1.3]\mittelpunkt{M}\sektor{0}{180}{}\bogen{0}{180}{b}\durchmesser{0}{d}\end{kreis}

\begin{teile}
\setcounter{teil}{3}
\teil Berechne die Länge des Bogens $b$. \quad \feld[cm]{b}
\teil Berechne die Länge des ganzen Randes der Halbkreisfläche (Bogen und Durchmesser). \\ Länge des Randes: \leerfeld[cm]
\teil Aus einem rechteckigen Blech wird ein Rohr mit dem Radius 4,3\,cm gebogen; die Kanten stoßen genau aneinander. Wie lang muss das Blech sein? Runde auf eine Stelle. (P10 2022)
\end{teile}
\end{aufgabe}

\begin{aufgabe}{Ich finde den Fehler beim Kreisumfang. Tom soll den Umfang eines Kreises mit dem Durchmesser $d = 14$\,cm berechnen. Er rechnet so:}
\rechnung{u &= 2 \cdot \pi \cdot r \\ u &= 2 \cdot \pi \cdot 14\,\text{cm} \\ u &\approx 88{,}0\,\text{cm}}
\begin{teile}
\teil Finde den Fehler, benenne ihn und korrigiere die Rechnung.
\teil Berechne den Umfang eines Kreises mit dem Durchmesser $d = 26$\,m.
\end{teile}
\end{aufgabe}

\begin{aufgabe}{Ich kann begründen, warum $u : d$ bei jedem Kreis gleich ist. Begründe ohne Rechnung.}
\begin{teile}
\teil Ein Kreis hat einen doppelt so großen Durchmesser wie ein anderer. Wie ändert sich der Umfang?
\teil Lena meint: „Bei einem großen Kreis ist $u : d$ größer als bei einem kleinen.“ Hat sie recht?
\teil Ein Rad mit 1\,m Durchmesser rollt einmal ganz herum. Ist der Weg länger oder kürzer als 3\,m?
\end{teile}
\end{aufgabe}

\begin{aufgabe}{Ich kann ausrechnen, wie weit ein Rad rollt. Das Vorderrad von Emils Fahrrad hat einen Durchmesser von 70\,cm.}
\begin{teile}
\teil Wie weit kommt das Rad bei einer Umdrehung? Gib das Ergebnis in cm an.
\teil Wie weit fährt Emil bei 500 Umdrehungen? Gib das Ergebnis in m an.
\teil Emils Schulweg ist 2,5\,km lang. Wie viele Umdrehungen macht das Vorderrad ungefähr?
\end{teile}
\end{aufgabe}
